\documentclass[11pt,letterpaper]{article}
% preamble.tex — Shared preamble for all documents
\usepackage{graphicx}
\usepackage{hyperref}
\usepackage{enumitem}
\usepackage{titlesec}
\usepackage{fancyhdr}
\usepackage{setspace}
\usepackage{parskip}
\usepackage{etoolbox}
\usepackage{xspace}
\usepackage{needspace}
\usepackage{changepage}
\usepackage{geometry}
\usepackage{ragged2e}
\usepackage{environ}

% Hyperref setup
\hypersetup{
  colorlinks=false,
  pdfborder={0 0 0},
  pdftitle={The Letter to the Father},
  pdfauthor={},
}

% Load project files
% colors.tex — Color palette for The Letter to the Father
\usepackage{xcolor}

% Primary palette — muted, period-appropriate, prestige
\definecolor{kafkablack}{HTML}{1A1A1A}
\definecolor{kafkagray}{HTML}{4A4A4A}
\definecolor{kafkalight}{HTML}{8A8A8A}
\definecolor{kafkacream}{HTML}{F5F0E8}
\definecolor{kafkasepia}{HTML}{704214}
\definecolor{kafkablood}{HTML}{6B1C23}
\definecolor{kafkaink}{HTML}{2C2C3A}

% Accent colors — for package materials
\definecolor{prestigegold}{HTML}{B8860B}
\definecolor{sectionrule}{HTML}{3A3A3A}
\definecolor{pullquote}{HTML}{5A4A3A}

% metadata.tex — Project metadata for The Letter to the Father
\newcommand{\filmtitle}{The Letter to the Father}
\newcommand{\filmtitleupper}{THE LETTER TO THE FATHER}
\newcommand{\filmsubtitle}{An Original Screenplay}
\newcommand{\filmauthor}{Based on the life and writings of Franz Kafka}
\newcommand{\filmyear}{2026}
\newcommand{\filmdraft}{Revised Draft -- Public Development Example}
\newcommand{\filmdate}{October 2026 -- Example Edition}
\newcommand{\filmlogline}{In the final years of his life, a frail and tormented Franz Kafka struggles to write, love, and exist within a world of family pressure, bureaucratic humiliation, illness, and spiritual dread --- while the stories he creates begin to mirror and consume the life he can no longer control.}
\newcommand{\filmgenre}{Prestige Drama / Literary Biopic}
\newcommand{\filmrating}{R}
\newcommand{\filmlead}{Adrien Brody}
\newcommand{\filmbudgetrange}{\$12--18M}
\newcommand{\filmruntime}{Approximately 110 minutes}


% Public example status travels with every downloadable document.
\AddToHook{shipout/foreground}{\put(40,-776){\makebox(0,0)[l]{\fontsize{8}{9}\selectfont\textit{Development example. All casting, financing and awards language is illustrative; no attachments confirmed.}}}}

% typography.tex — Typography settings
\usepackage{fontspec}
\usepackage{microtype}

% Prefer repo-local metric-compatible fonts so Tectonic does not depend on
% absolute host-system font paths. Fall back to system fonts only if the vendored
% font payload is unavailable.
\IfFileExists{../../../fonts/Tinos-Regular.ttf}{
  \setmainfont{Tinos-Regular.ttf}[
    Path=../../../fonts/,
    BoldFont=Tinos-Bold.ttf,
    ItalicFont=Tinos-Italic.ttf,
    BoldItalicFont=Tinos-BoldItalic.ttf,
  ]
}{
  \IfFontExistsTF{TeX Gyre Termes}{
    \def\projectmainfont{TeX Gyre Termes}
  }{
    \IfFontExistsTF{Times New Roman}{
      \def\projectmainfont{Times New Roman}
    }{
      \def\projectmainfont{Latin Modern Roman}
    }
  }
  \setmainfont{\projectmainfont}
}

\IfFileExists{../../../fonts/Arimo-Regular.ttf}{
  \setsansfont{Arimo-Regular.ttf}[
    Path=../../../fonts/,
    BoldFont=Arimo-Bold.ttf,
    ItalicFont=Arimo-Italic.ttf,
    BoldItalicFont=Arimo-BoldItalic.ttf,
  ]
}{
  \IfFontExistsTF{TeX Gyre Heros}{
    \def\projectsansfont{TeX Gyre Heros}
  }{
    \IfFontExistsTF{Helvetica}{
      \def\projectsansfont{Helvetica}
    }{
      \def\projectsansfont{Latin Modern Sans}
    }
  }
  \setsansfont{\projectsansfont}
}

\IfFileExists{../../../fonts/Cousine-Regular.ttf}{
  \setmonofont{Cousine-Regular.ttf}[
    Path=../../../fonts/,
    BoldFont=Cousine-Bold.ttf,
    ItalicFont=Cousine-Italic.ttf,
    BoldItalicFont=Cousine-BoldItalic.ttf,
    Scale=0.95,
  ]
  \def\screenplayfont{%
    \fontspec{Cousine-Regular.ttf}[%
      Path=../../../fonts/,
      BoldFont=Cousine-Bold.ttf,
      ItalicFont=Cousine-Italic.ttf,
      BoldItalicFont=Cousine-BoldItalic.ttf,
      Scale=1.0,
    ]%
  }
}{
  \IfFontExistsTF{Courier New}{
    \def\projectmonofont{Courier New}
  }{
    \IfFontExistsTF{Courier}{
      \def\projectmonofont{Courier}
    }{
      \def\projectmonofont{Latin Modern Mono}
    }
  }
  \setmonofont{\projectmonofont}[Scale=0.95]
  \def\screenplayfont{\fontspec{\projectmonofont}[Scale=1.0]}
}

% Microtype settings
\microtypesetup{
  protrusion=true,
  expansion=false,
}

% studio_package_macros.tex — Formatting macros for the studio package

% ============================================================
% SECTION STYLING
% ============================================================
\titleformat{\section}
  {\Large\bfseries\color{kafkablack}\sffamily}
  {}
  {0em}
  {}
  [\vspace{2pt}\textcolor{sectionrule}{\rule{\textwidth}{0.5pt}}\vspace{6pt}]

\titleformat{\subsection}
  {\large\bfseries\color{kafkagray}\sffamily}
  {}
  {0em}
  {}

\titleformat{\subsubsection}
  {\normalsize\bfseries\color{kafkagray}\sffamily}
  {}
  {0em}
  {}

% ============================================================
% PULL QUOTE
% ============================================================
% Usage: \pullquote{Text here}
\NewEnviron{pullquote}{%
  \vspace{8pt}%
  \begin{adjustwidth}{0.5in}{0.5in}%
    \color{pullquote}\itshape\large%
    \BODY%
  \end{adjustwidth}%
  \vspace{8pt}%
}

% ============================================================
% HIGHLIGHT BOX
% ============================================================
\newcommand{\highlightbox}[1]{%
  \par\vspace{6pt}%
  \noindent\colorbox{kafkacream}{%
    \parbox{\dimexpr\textwidth-2\fboxsep}{%
      \vspace{4pt}#1\vspace{4pt}%
    }%
  }%
  \par\vspace{6pt}%
}

% ============================================================
% TARGET ENTRY (for director/producer lists)
% ============================================================
% Usage: \target{Name}{Why they fit}{Tier}
\newcommand{\target}[3]{%
  \needspace{4\baselineskip}%
  \subsubsection*{#1 \hfill \textnormal{\small\textit{#3}}}%
  #2\par\vspace{6pt}%
}

% ============================================================
% SCENE CARD (for showcase section)
% ============================================================
% Usage: \scenecard{Title}{Description}
\newcommand{\scenecard}[2]{%
  \needspace{5\baselineskip}%
  \vspace{4pt}%
  \noindent\textbf{\sffamily #1}\\[3pt]%
  #2\par\vspace{8pt}%
  \textcolor{sectionrule}{\rule{0.3\textwidth}{0.25pt}}%
  \vspace{8pt}%
}

% ============================================================
% TAGLINE
% ============================================================
\newcommand{\tagline}[1]{%
  \vspace{4pt}%
  \begin{center}%
    {\large\itshape\color{kafkagray} ``#1''}%
  \end{center}%
  \vspace{4pt}%
}

% ============================================================
% COMP FILM ENTRY
% ============================================================
% Usage: \compfilm{Title (Year)}{Relevance}
\newcommand{\compfilm}[2]{%
  \noindent\textbf{#1}\\%
  {\small #2}\par\vspace{6pt}%
}

% ============================================================
% RISK ENTRY
% ============================================================
% Usage: \risk{Risk}{Mitigation}
\newcommand{\risk}[2]{%
  \noindent\textbf{Risk:} #1\\%
  \textbf{Mitigation:} #2\par\vspace{8pt}%
}

% ============================================================
% PACKAGE GEOMETRY
% ============================================================
\newcommand{\packagegeometry}{%
  \newgeometry{
    left=1in,
    right=1in,
    top=1in,
    bottom=1in,
  }%
}

\geometry{left=0.75in, right=0.75in, top=0.75in, bottom=0.75in}
\pagestyle{empty}
\setlength{\parindent}{0pt}
\setlength{\parskip}{5pt}

\begin{document}

\begin{center}
{\Large\bfseries\sffamily \filmtitleupper}\\[3pt]
{\normalsize\sffamily\color{kafkagray} Director Attachment Brief}
\end{center}
\vspace{2pt}
\textcolor{sectionrule}{\rule{\textwidth}{0.5pt}}
\vspace{4pt}

\textbf{The Proposition}\\
A prestige biographical drama with unusually clear directing engines for this kind of material. The office is a machine. The father-letter is a title object that keeps returning. Berlin is the first room in which Kafka is briefly happy, which makes the collapse hurt more. The film does what most biopics cannot: it makes Kafka's real life take on the emotional and structural logic of a Kafka story within rigorous prestige realism.

\vspace{2pt}

\textbf{Formal Language}\\
The film operates through six emotional movements (not chronological acts) anchored in Kafka's final years, with earlier material arriving as contamination rather than explanatory flashback. The visual grammar is built around recurring systems: corridors subtly elongate, social scenes block as tribunals, administrative sounds leak into domestic spaces, and physical space compresses as Kafka's body fails. Crucially, the undelivered Letter to His Father is not just a mid-film centerpiece but a contaminant that re-poisons the Berlin and Kierling material. The method is disciplined subjectivity, not dream-sequence surrealism: controlled destabilization inside a world that still feels physically real.

\vspace{2pt}

\textbf{Visual Opportunities}\\
\textit{Set Pieces}: the elevator as metaphysical exclusion; the engagement dinner as tribunal; the returned Letter entering the Berlin apartment like a curse; the ``burn everything'' argument with Brod; the final room of handwritten slips. --- \textit{Space}: The Insurance Institute as existential geometry; the sanatorium as bureaucracy of the body. --- \textit{Light}: Amber Prague interiors yielding to cold Berlin gray yielding to the cruel spring brightness of the death room. --- \textit{Sound}: Stamping, filing, coughing, paper, waiting. The institutional soundscape of a life lived inside systems.

\vspace{2pt}

\textbf{The Actor}\\
Adrien Brody. Post-\textit{Brutalist}. At the peak of his powers and his credibility. The physical and temperamental match with Kafka is precise. Six major performance scenes, multiple passages of minimal dialogue, a visible deterioration arc, and a role that moves from public competence to private collapse without ever losing intelligence. A genuine creative partnership opportunity.

\vspace{2pt}

\textbf{Festival and Awards}\\
Venice Competition premiere target. Best Actor as primary category, with Director, Picture, Screenplay, and Cinematography all strongly competitive. This is a film designed for the highest tier of international festival and awards recognition.

\vspace{2pt}

\textbf{Production}\\
Budget: \filmbudgetrange. Czech Republic production base with tax incentives. 50--55 day shoot. Interior-heavy, period-specific, formally demanding. Requires a director who understands that the most powerful cinema happens in the space between what is said and what is seen.

\vspace{2pt}

\textbf{What This Film Needs from a Director}\\
Formal control. Patience. The ability to make institutions cinematic without inflating them into gimmickry. An understanding that the Berlin movement must feel genuinely warm or the whole film collapses into monochrome suffering. Comfort with silence, shame, and performance carried in the face. The willingness to let a scene sit in discomfort. The ability to collaborate with an actor of Brody's caliber at the deepest level. And the ambition to make a biographical drama that transcends the form.

\vspace{4pt}
\textcolor{sectionrule}{\rule{\textwidth}{0.5pt}}
\vspace{2pt}
\begin{center}
{\small\color{kafkalight}\textit{A biographical drama with the shape of a Kafka story and the discipline of prestige realism.}}
\end{center}

\end{document}
